% ============================================================
%  Geometría Analítica · LM · Universidad de Sonora
%  Semana 2 — Vectores, la recta en el plano y el producto punto
%  Semana CON laboratorio (el laboratorio es quincenal)
%  Compilar en Overleaf con pdfLaTeX
% ============================================================
\documentclass[aspectratio=169,10pt]{beamer}
\usepackage[utf8]{inputenc}
\usepackage[T1]{fontenc}
\usepackage[spanish,es-nodecimaldot,es-noquoting]{babel}
\usepackage{lmodern}
\usepackage{amsmath,amssymb}
\usepackage{booktabs}
\usepackage{tikz}
\usetikzlibrary{arrows.meta,calc}
\usepackage{graphicx}
\graphicspath{{figuras/}}
% ---------- Tema y colores ----------
\usetheme{default}
\setbeamertemplate{navigation symbols}{}
\definecolor{turquesaGA}{HTML}{0E7C8C}
\definecolor{turquesaClaro}{HTML}{E2F2F4}
\definecolor{grisT}{HTML}{5F5E5A}
\definecolor{negroP}{HTML}{2B2B2B}
\definecolor{terracota}{HTML}{C1502E}
\definecolor{dorado}{HTML}{B07C10}
\definecolor{verdeOk}{HTML}{1D9E75}
\definecolor{rojoAc}{HTML}{C0392B}
\setbeamercolor{structure}{fg=turquesaGA}
\setbeamercolor{frametitle}{fg=white,bg=turquesaGA}
\setbeamercolor{title}{fg=turquesaGA}
\setbeamerfont{frametitle}{size=\large,series=\bfseries}
\setbeamertemplate{frametitle}{%
  \nointerlineskip
  \begin{beamercolorbox}[wd=\paperwidth,ht=2.6ex,dp=1.2ex,leftskip=0.3cm]{frametitle}
    \usebeamerfont{frametitle}\insertframetitle
  \end{beamercolorbox}}
\setbeamertemplate{footline}{%
  \hfill\textcolor{grisT}{\tiny Geometría Analítica · LM · Unison \quad
  \insertframenumber/\inserttotalframenumber}\hspace{0.3cm}\vspace{0.15cm}}
\setbeamertemplate{itemize item}{\textcolor{turquesaGA}{\raisebox{0.15ex}{\tiny$\blacksquare$}}}
\setbeamertemplate{itemize subitem}{\textcolor{grisT}{\raisebox{0.15ex}{\tiny$\bullet$}}}
% ---------- Caja auxiliar para destacar ideas ----------
\newcommand{\destaca}[1]{%
  \par\vspace{3pt}\noindent
  \colorbox{turquesaClaro}{%
    \parbox{\dimexpr\linewidth-2\fboxsep\relax}{\vspace{3pt}#1\vspace{3pt}}}%
  \par\vspace{3pt}}
% ---------- Hojas de separación por sesión ----------
%  Las sesiones no llevan número: se identifican por su título.
\newcommand{\hojaSesion}[2]{%
  \vfill
  \begin{beamercolorbox}[wd=\textwidth,sep=16pt,center]{frametitle}
    {\Large\bfseries #1}\\[11pt]
    \begin{minipage}{0.82\textwidth}\centering\small #2\end{minipage}
  \end{beamercolorbox}
  \vfill}
% ---------- Abreviaturas ----------
\newcommand{\R}{\mathbb{R}}
\newcommand{\pp}[2]{\langle #1,#2\rangle}
\newcommand{\norma}[1]{\lVert #1\rVert}
% ============================================================
\title{\LARGE\textbf{Geometría Analítica}}
\subtitle{Semana 2 · Vectores, la recta en el plano y el producto punto}
\author{Prof. Rosalía Gpe. Hernández Amador}
\institute{Departamento de Matemáticas · Universidad de Sonora}
\date{Semestre 2026-2}
\begin{document}
% ============================================================
\begin{frame}[plain]
  \titlepage
\end{frame}
% ============================================================
\begin{frame}{Ruta de la semana}
  \centering
  {\small
  \begin{tabular}{@{}l@{}}
    \toprule
    \textbf{Las cinco sesiones de la semana, en orden} \\
    \midrule
    Puntos y vectores: un objeto nuevo y sus operaciones \\
    La recta descrita con un vector: forma paramétrica y pendiente \\
    El vector normal y la ecuación general; las formas de la recta \\
    Taller de las dos direcciones; la mediatriz \\
    El producto punto: definición, propiedades y el teorema del ángulo \\
    \midrule
    Laboratorio: Construcción 1 en GeoGebra (esta semana hay laboratorio) \\
    \bottomrule
  \end{tabular}}

  \vspace{0.4cm}
  \raggedright
  \destaca{La semana pasada describimos rectas donde una coordenada queda
  fija ($x=3$, $y=-2$). Esta semana construimos el \textbf{vector}, la
  herramienta que describe \textbf{cualquier} recta, y al cerrar la
  semana definimos el producto punto.}
\end{frame}
% ============================================================
\section{Puntos y vectores}
% ---------- Portada de la Unidad 2 ----------
%  Misma página ligera de apertura de unidad que en la semana 1.
\begin{frame}[plain]
  \begin{columns}[c]
    \begin{column}{0.42\textwidth}
      \centering
      \begin{tikzpicture}[scale=0.6,>=Stealth]
        \draw[very thin,color=gray!30] (-0.9,-1.2) grid (4.6,3.1);
        \draw[->] (-0.9,0) -- (4.8,0);
        \draw[->] (0,-1.2) -- (0,3.3);
        \draw[turquesaGA,very thick] (-0.7,-0.85) -- (4.4,1.7);
        \draw[terracota,very thick,->] (1,0) -- (3,1)
          node[midway,above,sloped,font=\scriptsize]{$d$};
      \end{tikzpicture}
    \end{column}
    \begin{column}{0.52\textwidth}
      \raggedright
      {\small\color{grisT} SEMANA 2}\\[8pt]
      {\color{dorado}\rule{0.55\linewidth}{1.6pt}}\\[12pt]
      {\LARGE\bfseries\color{turquesaGA} Unidad 2}\\[8pt]
      {\large La recta en el plano}\\[14pt]
      {\small\color{grisT}\itshape Todas las rectas, descritas con un
      mismo lenguaje: el de los vectores}
    \end{column}
  \end{columns}
\end{frame}
% ---------- hoja de sesión ----------
\begin{frame}[plain]
  \hojaSesion{Puntos y vectores}{%
  El vector como desplazamiento. Qué operaciones entre puntos y
  vectores tienen sentido geométrico y cuáles no, por qué la suma de
  dos puntos no está definida, y cómo la expresión $P+t\,(Q-P)$
  reescribe el punto medio y describe la recta completa.}
\end{frame}
\begin{frame}{Pregunta pendiente}
  Cerramos la semana 1 con esta pregunta. Si $xy>0$ y $x+y<0$,
  ¿en qué cuadrante se encuentra el punto $(x,y)$?

  \vspace{0.3cm}
  \begin{itemize}
    \item $xy>0$: los signos de $x$ e $y$ son \textbf{iguales},
      así que el punto está en el cuadrante I o en el III.
    \item $x+y<0$: la suma de dos positivos no puede ser negativa,
      así que se descarta el cuadrante I.
  \end{itemize}

  \vspace{0.3cm}
  \destaca{El punto está en el \textbf{cuadrante III}. Nótese el método:
  cada condición \textbf{recorta} el conjunto de posibilidades, y la
  respuesta es lo que sobrevive a todos los recortes. Así razonamos con
  los lugares geométricos durante todo el curso.}
\end{frame}
% ============================================================
\begin{frame}{Punto y vector}
  Un \textcolor{turquesaGA}{\textbf{punto}} es un \textit{lugar}.
  Un \textcolor{terracota}{\textbf{vector}} es un \textit{desplazamiento}.
  Ambos se escriben como parejas de números, y de ahí viene la posible
  confusión. Las operaciones con sentido geométrico son:

  \vspace{0.3cm}
  \begin{center}
  {\small
  \begin{tabular}{@{}lll@{}}
    \toprule
    \textbf{Expresión} & \textbf{Naturaleza} & \textbf{Interpretación} \\
    \midrule
    $Q-P$ & \textcolor{terracota}{vector} & el desplazamiento de $P$ a $Q$ \\
    $P+v$ & \textcolor{turquesaGA}{punto} & $P$ desplazado según $v$ \\
    $u+v$, \ $\lambda u$ & \textcolor{terracota}{vectores} & componer y reescalar desplazamientos \\
    $P+Q$ & --- & \textcolor{rojoAc}{\textbf{no está definido}} \\
    \bottomrule
  \end{tabular}}
  \end{center}

  \vspace{0.3cm}
  \destaca{Regla de lectura: punto $-$ punto $=$ vector; \ punto $+$
  vector $=$ punto; \ punto $+$ punto no significa nada. Antes de operar
  con parejas de números, conviene preguntarse qué \textbf{es} cada una.}
\end{frame}
% ============================================================
\begin{frame}{¿Por qué $P+Q$ no está definido?}
  Pensemos en un mapa de Sonora:

  \vspace{0.2cm}
  \begin{itemize}
    \item \textit{El desplazamiento de Hermosillo a Guaymas} tiene un
      significado claro, independiente de dónde coloquemos el origen
      del mapa.

      \vspace{0.2cm}
    \item \textit{Hermosillo más Guaymas} no corresponde a ninguna
      ciudad.
  \end{itemize}

  \vspace{0.4cm}
  Si de cualquier modo sumamos las coordenadas, el resultado
  \textbf{depende de dónde se colocó el origen}: al mover el origen del
  mapa, la ``suma'' cae en otro sitio.

  \vspace{0.3cm}
  \destaca{Eso revela que $P+Q$ no es un objeto geométrico. La geometría
  no puede depender de una elección arbitraria (la posición del origen).
  En cambio, $Q-P$ \textbf{no cambia} al mover el origen, y por eso el
  desplazamiento sí es un objeto legítimo.}
\end{frame}
% ============================================================
\begin{frame}{El punto medio}
  Del bachillerato conocemos la fórmula
  $M=\left(\frac{x_1+x_2}{2},\frac{y_1+y_2}{2}\right)$.
  Pero acabamos de acordar que la suma de puntos no está definida.
  La escritura correcta es
  \[
    M \;=\; P+\tfrac{1}{2}\,(Q-P).
  \]
  Partimos de $P$ y avanzamos \textit{la mitad del desplazamiento}
  hacia $Q$. Punto más vector es punto, así que cada paso es legítimo,
  y las coordenadas que resultan son las de siempre.

  \vspace{0.3cm}
  \destaca{\textbf{Con $t$ arbitrario}, la expresión $P+t\,(Q-P)$ produce

  \begin{itemize}
    \item para $t=0$, el punto $P$; para $t=1$, el punto $Q$
    \item para $t=\tfrac12$, el punto medio
    \item al recorrer $t$ todos los reales, \textbf{toda la recta que
      pasa por $P$ y $Q$}
  \end{itemize}
  Con esta observación empezamos la siguiente sesión.}
\end{frame}
% ============================================================
\begin{frame}{El applet de la semana}
  \begin{columns}[T]
    \begin{column}{0.54\textwidth}
      En el applet \textit{Puntos y vectores} del sitio del curso,
      antes de mover la figura, escriban su predicción:

      \vspace{0.2cm}
      \begin{enumerate}
        \item Si $P$ y $Q$ se trasladan \textbf{juntos},
          ¿cambia el vector $Q-P$?
        \item ¿Y si se desplaza únicamente $Q$?
      \end{enumerate}

      \vspace{0.3cm}
      \destaca{\small Conclusión esperada: un desplazamiento no depende
      de dónde inicia, sino únicamente de su \textbf{magnitud} y su
      \textbf{dirección}. Esa invariancia distingue al vector del
      punto.}
    \end{column}
    \begin{column}{0.42\textwidth}
      \begin{center}
      \begin{tikzpicture}[scale=.6,>=Stealth]
        \draw[very thin,color=gray!30] (-2.4,-1.9) grid (3.4,2.9);
        \draw[->] (-2.4,0) -- (3.4,0);
        \draw[->] (0,-1.9) -- (0,2.9);
        \coordinate (P) at (-1.5,-1);
        \coordinate (Q) at (2.2,1.6);
        \draw[terracota,very thick,->] (P) -- (Q)
          node[midway,above,sloped]{\small $Q-P$};
        \fill[turquesaGA] (P) circle(2.4pt) node[below left]{$P$};
        \fill[dorado] (Q) circle(2.4pt) node[above right]{$Q$};
        \coordinate (P2) at (-0.6,0.1);
        \coordinate (Q2) at (3.1,2.7);
        \draw[terracota!40,very thick,->] (P2) -- (Q2);
        \fill[turquesaGA!40] (P2) circle(2.4pt);
        \fill[dorado!40] (Q2) circle(2.4pt);
      \end{tikzpicture}
      \end{center}
    \end{column}
  \end{columns}
\end{frame}
% ============================================================
\section{La recta descrita con un vector}
% ---------- hoja de sesión ----------
\begin{frame}[plain]
  \hojaSesion{La recta descrita con un vector}{%
  Una recta queda determinada por un punto de paso y una dirección:
  la forma paramétrica $X=P+t\,d$ y el vector director. El puente con
  el lenguaje del bachillerato ---pendiente e inclinación--- y el caso
  donde ese lenguaje falla y el vectorial no: las rectas verticales.}
\end{frame}
\begin{frame}{¿Qué determina una recta?}
  Una recta queda determinada por
  \begin{center}
    {\large un \textbf{punto de paso} $P$ \quad y \quad
     una \textbf{dirección} $d$.}
  \end{center}

  \vspace{0.3cm}
  \destaca{\textbf{Forma paramétrica (o vectorial).}
  \[
    \ell:\quad X = P + t\,d, \qquad t\in\R .
  \]
  El vector $d$ se llama \textbf{vector director} de la recta. Cada
  valor del parámetro $t$ produce un punto; al recorrer $t$ todos los
  reales, recorremos toda la recta.}

  \vspace{0.3cm}
  Nótese que es la expresión $P+t\,(Q-P)$ de la sesión anterior, con
  $d=Q-P$. Conocer dos puntos de una recta nos da de inmediato un
  vector director.
\end{frame}
% ============================================================
\begin{frame}{Ejemplos}
  \textbf{Ejemplo 1.} La recta pasa por $P=(1,2)$ con director
  $d=(3,1)$.
  \begin{enumerate}
    \item Determinar los puntos correspondientes a $t=0,\ 1,\ 2,\ -1$.
    \item ¿Pertenece el punto $(10,5)$ a la recta? ¿Y el $(7,5)$?
  \end{enumerate}

  \vspace{0.4cm}
  \destaca{\small \textit{Solución.} 1.~$(1,2)$, $(4,3)$, $(7,4)$,
  $(-2,1)$.\quad 2.~El punto $(10,5)$ sí pertenece: $t=3$ satisface
  ambas coordenadas. El punto $(7,5)$ no: la abscisa exige $t=2$ y la
  ordenada exige $t=3$, y \textbf{el parámetro debe ser el mismo en
  ambas coordenadas}. Este es el criterio de pertenencia que usaremos
  en el cuestionario de la semana.}
\end{frame}
% ============================================================
\begin{frame}{De la dirección a la pendiente}
  \begin{columns}[T]
    \begin{column}{0.52\textwidth}
      Con director $d=(d_1,d_2)$ y $d_1\neq 0$, por cada $d_1$ unidades
      de avance en $x$ la ordenada cambia $d_2$. La razón
      \[
        m=\frac{d_2}{d_1}
      \]
      es la \textbf{pendiente}. Si la recta pasa por $P=(x_1,y_1)$ y
      $Q=(x_2,y_2)$:
      \[
        m=\frac{y_2-y_1}{x_2-x_1}.
      \]
      La \textbf{inclinación} es el ángulo $\alpha\in[0,\pi)$ que la
      recta forma con el semieje $X$ positivo, y $m=\tan\alpha$.
    \end{column}
    \begin{column}{0.44\textwidth}
      \begin{center}
      \begin{tikzpicture}[scale=.64,>=Stealth]
        \draw[very thin,color=gray!30] (-0.9,-0.9) grid (4.6,3.1);
        \draw[->] (-0.9,0) -- (4.6,0) node[below]{$x$};
        \draw[->] (0,-0.9) -- (0,3.1) node[left]{$y$};
        \draw[turquesaGA,very thick] (-0.7,-0.3) -- (4.4,2.25);
        \coordinate (A) at (1,0.55);
        \coordinate (B) at (3,1.55);
        \draw[terracota,very thick,->] (A) -- (B)
          node[midway,above,sloped]{\small $d$};
        \draw[dashed,dorado,thick] (A) -- (3,0.55)
          node[midway,below]{\small $d_1$};
        \draw[dashed,dorado,thick] (3,0.55) -- (B)
          node[midway,right]{\small $d_2$};
      \end{tikzpicture}
      \end{center}
    \end{column}
  \end{columns}

  \vspace{0.05cm}
  \destaca{\small Las rectas \textbf{verticales} no tienen pendiente
  (con $d_1=0$ la razón no está definida), pero sí tienen vector
  director, $(0,1)$. El lenguaje vectorial cubre los casos donde el de
  pendientes falla; por eso será el lenguaje principal del curso.}
\end{frame}
% ============================================================
\begin{frame}{Paralelismo en ambos lenguajes}
  \centering
  {\small
  \begin{tabular}{@{}lcc@{}}
    \toprule
    & \textbf{con pendientes} & \textbf{con directores} \\
    \midrule
    $\ell_1 \parallel \ell_2$ & $m_1=m_2$ & $d^{(1)}=\lambda\, d^{(2)}$ \\
    $\ell_1 \perp \ell_2$ & $m_1\, m_2=-1$ & (próxima sesión) \\
    \bottomrule
  \end{tabular}}
  \par\vspace{0.5cm}
  \raggedright
  Obsérvese que multiplicar el director por $\lambda\neq 0$ no cambia
  la recta que describe. Por ejemplo, $(4,-6)$ y $(2,-3)$ dirigen la
  misma familia de rectas paralelas. La perpendicularidad con vectores
  la construiremos en la siguiente sesión, a partir de un giro de
  $90^\circ$.

  \vspace{0.3cm}
  \destaca{\textit{Ejercicio.} Determinar un vector director de la
  recta $X=(2,-1)+t\,(4,-6)$ con componentes enteras de valor absoluto
  mínimo, y calcular su pendiente.
  \\[3pt]{\small \textit{Solución:} $(2,-3)$, con
  $m=-\tfrac{3}{2}$.}}
\end{frame}
% ============================================================
\section{El vector normal y la ecuación general}
% ---------- hoja de sesión ----------
\begin{frame}[plain]
  \hojaSesion{El vector normal y la ecuación general}{%
  El giro de noventa grados produce el vector normal. En la
  verificación aparece una suma de productos que se anula y que pronto
  recibirá nombre propio. La ecuación general $Ax+By=C$, cuyos
  coeficientes son las coordenadas de un vector normal, y las formas de
  la ecuación de la recta.}
\end{frame}
\begin{frame}{El vector perpendicular}
  \begin{columns}[T]
    \begin{column}{0.52\textwidth}
      Dado un vector $d=(d_1,d_2)$, construimos
      \[
        n=(-d_2,\ d_1)
      \]
      (intercambiamos las componentes y cambiamos el signo de una de
      ellas).

      \vspace{0.2cm}

      \textbf{Afirmación:} $n$ es $d$ girado $90^\circ$.
      Verificación con pendientes:
      \[
        \frac{d_2}{d_1}\cdot\frac{d_1}{-d_2}=-1 .
      \]
    \end{column}
    \begin{column}{0.44\textwidth}
      \begin{center}
      \begin{tikzpicture}[scale=.72,>=Stealth]
        \draw[very thin,color=gray!30] (-2.6,-0.9) grid (3.1,3.1);
        \draw[->] (-2.6,0) -- (3.1,0);
        \draw[->] (0,-0.9) -- (0,3.1);
        \draw[terracota,very thick,->] (0,0) -- (2.5,1)
          node[right]{\small $d=(d_1,d_2)$};
        \draw[turquesaGA,very thick,->] (0,0) -- (-1,2.5)
          node[above]{\small $n=(-d_2,d_1)$};
        \draw (0.35,0.14) arc (22:112:0.38);
        \node[font=\scriptsize] at (0.16,0.52) {$90^\circ$};
      \end{tikzpicture}
      \end{center}
    \end{column}
  \end{columns}

  \vspace{0.2cm}
  \destaca{Obsérvese además la expresión
  $d_1\cdot(-d_2)+d_2\cdot d_1=0$, una \textbf{suma de productos de
  componentes} que se anula justo cuando los vectores son
  perpendiculares. Esta expresión \textbf{recibirá nombre propio al cerrar la
  semana}; con ella mediremos todos los ángulos y todas las
  distancias del curso.}
\end{frame}
% ============================================================
\begin{frame}{El vector normal y la ecuación general}
  Una recta también queda determinada por un punto $P=(x_0,y_0)$ y un
  vector \textbf{perpendicular} a ella: su \textbf{vector normal}
  $n=(A,B)$. La condición de pertenencia se escribe
  \[
    Ax+By=C, \qquad\text{con } C=Ax_0+By_0
  \]
  (obtenemos $C$ exigiendo que $P$ satisfaga la ecuación).
  A esta expresión la llamamos \textbf{ecuación general} de la recta.

  \vspace{0.3cm}
  \destaca{\textbf{Observación importante.} Los coeficientes $A$ y $B$
  de la ecuación general \textbf{son las coordenadas de un vector
  normal} a la recta. En la ecuación $3x+y=5$, el vector $(3,1)$ es
  perpendicular a la recta: esa información está a la vista, sin
  calcular nada. Saber \textbf{leer} una ecuación es tan importante
  como saber obtenerla.}
\end{frame}
% ============================================================
\begin{frame}{Formas de la ecuación de la recta}
  \centering
  {\small
  \begin{tabular}{@{}lll@{}}
    \toprule
    \textbf{Forma} & \textbf{Expresión} & \textbf{Dato que permite leer} \\
    \midrule
    Paramétrica & $X=P+t\,d$ & un punto y el director; válida en $\R^3$ \\
    Punto--pendiente & $y-y_0=m(x-x_0)$ & un punto y la pendiente \\
    Pendiente--ordenada & $y=mx+b$ & la pendiente y el corte con el eje $Y$ \\
    General & $Ax+By=C$ & el normal $(A,B)$; incluye verticales \\
    Simétrica & $\dfrac{x}{a}+\dfrac{y}{b}=1$ & los cortes $(a,0)$ y $(0,b)$ \\
    \bottomrule
  \end{tabular}}

  \vspace{0.5cm}
  \raggedright
  \destaca{La habilidad que practicamos \textbf{no} consiste en
  memorizar las cinco formas, sino en \textbf{traducir} entre ellas.
  Cada una exhibe un dato distinto de la misma recta, y el problema
  dicta cuál conviene usar.}
\end{frame}
% ============================================================
\begin{frame}{Ejemplos}
  \textbf{Ejemplo 2.} La recta pasa por $(2,1)$ y $(6,3)$.
  \begin{enumerate}
    \item Determinar su forma paramétrica y su pendiente.
    \item Escribir la forma punto--pendiente y llevarla a $y=mx+b$.
    \item Escribir la forma general y leer su vector normal.
    \item Comprobar que ese normal es perpendicular al director
      (la suma de productos correspondiente debe anularse).
  \end{enumerate}

  \vspace{0.4cm}
  \destaca{\small \textit{Solución.} $X=(2,1)+t\,(4,2)$;
  $m=\tfrac12$; \ $y-1=\tfrac12(x-2)$, es decir $y=\tfrac12 x$; \
  forma general $x-2y=0$, con normal $(1,-2)$; \ y
  $4\cdot 1+2\cdot(-2)=0$, como esperábamos. Nótese que con la
  verificación del inciso 4 \textbf{detectamos errores propios} antes
  de que los detecte el examen.}
\end{frame}
% ============================================================
\section{Taller y laboratorio}
% ---------- hoja de sesión ----------
\begin{frame}[plain]
  \hojaSesion{Taller: las dos direcciones}{%
  Práctica de las dos traducciones por separado, de la ecuación a la
  gráfica y de la descripción a la ecuación. La mediatriz como lugar
  geométrico, con sus dos comprobaciones, y la Construcción~1 en
  GeoGebra con su criterio de evaluación, la prueba de arrastre.}
\end{frame}
\begin{frame}{De la ecuación a la gráfica}
  Para cada recta, identificamos la forma, leemos el dato que está a
  la vista y describimos un procedimiento para trazarla.

  \vspace{0.2cm}
  \begin{center}
  {\small
  \begin{tabular}{@{}lcc@{}}
    \toprule
    \textbf{Recta} & \textbf{¿Qué forma es?} & \textbf{¿Qué dato está a la vista?} \\
    \midrule
    $y=2x-3$ & & \\[8pt]
    $3x+4y=12$ & & \\[8pt]
    $X=(0,1)+t\,(1,-1)$ & & \\[8pt]
    $x=5$ & & \\
    \bottomrule
  \end{tabular}}
  \end{center}

  \vspace{0.4cm}
  {\footnotesize La tabla la llenamos en voz alta durante el taller.
  Obsérvese el último renglón: esa recta no tiene pendiente, pero sí
  tiene vector director.}
  % Guía del taller (las soluciones completas van en las notas):
  %  (a) pendiente-ordenada; corte (0,-3) y m=2.
  %  (b) general; normal (3,4), cortes (4,0) y (0,3).
  %  (c) paramétrica; punto de paso (0,1) y director (1,-1).
  %  (d) vertical; sin pendiente, director (0,1).
\end{frame}
% ============================================================
\begin{frame}{De la descripción a la ecuación}
  \destaca{\textbf{Problema.} Determinar el lugar geométrico de los
  puntos que \textbf{equidistan} de $A=(-2,1)$ y $B=(4,3)$.}

  \vspace{0.2cm}
  Antes de calcular, escribimos la predicción. ¿Qué figura esperamos,
  y por qué?

  \vspace{0.2cm}

  {\small \textit{Solución.} La condición es $d(P,A)=d(P,B)$ con
  $P=(x,y)$. Al elevar al cuadrado (operación reversible aquí, pues
  ambos lados son no negativos) y cancelar $x^2$, $y^2$:
  \[
    \boxed{\,3x+y=5\,}
  \]
  Una \textbf{recta}, la \textbf{mediatriz} del segmento $AB$.}

  \vspace{0.2cm}
  \destaca{\small Dos comprobaciones inmediatas: el punto medio $(1,2)$
  satisface la ecuación, y el normal $(3,1)$ es paralelo a
  $B-A=(6,2)$, es decir, la mediatriz es perpendicular al segmento,
  como debía ser. Verificar por una segunda vía es uno de nuestros dos
  acuerdos de trabajo.}
\end{frame}
% ============================================================
\begin{frame}{Laboratorio: Construcción 1}
  \textbf{Construcción 1} (actividad GeoGebra en Moodle):
  \begin{enumerate}
    \item Dos puntos libres $P$ y $Q$.
    \item La recta que pasa por ellos.
    \item El vector director $Q-P$ y, construido a partir de él,
      un vector normal.
  \end{enumerate}

  \vspace{0.4cm}
  \destaca{\textbf{Criterio de evaluación: la prueba de arrastre.}
  Al arrastrar $P$ o $Q$, todas las propiedades deben conservarse: la
  recta sigue pasando por ambos puntos y el normal permanece
  perpendicular. Una figura que se deshace al arrastrarla estaba
  \textit{dibujada}; una que resiste estaba \textit{construida}.}

  \vspace{0.3cm}
  {\footnotesize La construcción la empezamos en el laboratorio; la
  terminan en casa y la entregan en Moodle el domingo.}
\end{frame}
% ============================================================
\section{El producto punto}
% ---------- hoja de sesión ----------
\begin{frame}[plain]
  \hojaSesion{El producto punto}{%
  La suma de productos que apareció con el vector normal recibe
  nombre. Definición y
  propiedades del producto punto, el teorema que lo conecta con el
  ángulo entre vectores, y la demostración que quedó pendiente en la
  semana 1, la desigualdad del triángulo.}
\end{frame}
\begin{frame}{Definición del producto punto}
  Al construir el vector normal, giramos $d=(d_1,d_2)$ hacia
  $n=(-d_2,d_1)$ y apareció una expresión cuyo valor fue cero:
  \[
    d_1\cdot(-d_2)+d_2\cdot d_1=0 .
  \]
  Anunciamos que esa suma de productos recibiría nombre propio.
  Hoy cumplimos.

  \vspace{0.2cm}
  \destaca{\textbf{Definición (producto punto).} Para $u=(u_1,u_2)$ y
  $v=(v_1,v_2)$:
  \[
    \pp{u}{v} \;=\; u_1v_1+u_2v_2 .
  \]
  En el espacio, la definición es análoga, con tres sumandos.}

  \vspace{0.2cm}
  Obsérvese que la operación recibe \textbf{dos vectores} y produce
  \textbf{un número}, no un vector. Eso la hace útil para
  \textbf{medir}.
\end{frame}
% ============================================================
\begin{frame}{Propiedades básicas}
  Se verifican directamente de la definición:

  \vspace{0.2cm}
  \begin{enumerate}
    \item \textbf{Simetría:} $\pp{u}{v}=\pp{v}{u}$.
    \item \textbf{Linealidad:} $\pp{\lambda u+w}{v}
      =\lambda\pp{u}{v}+\pp{w}{v}$.
    \item \textbf{Positividad:} $\pp{u}{u}=u_1^2+u_2^2=\norma{u}^2$.
  \end{enumerate}

  \vspace{0.4cm}
  La tercera propiedad conecta con lo ya conocido:
  \[
    \norma{u}=\sqrt{\pp{u}{u}}
    \qquad\text{y por tanto}\qquad
    d(P,Q)=\norma{Q-P}.
  \]

  \vspace{0.1cm}
  \destaca{La distancia de la semana 1 ya era un producto punto. Lo
  nuevo no es la fórmula sino la operación con sus propiedades; de
  ellas saldrán los ángulos y las distancias que faltan.}
\end{frame}
% ============================================================
\begin{frame}{El teorema del ángulo}
  \destaca{\textbf{Teorema (el producto punto mide ángulos).}
  Si $u,v\neq 0$ y $\theta\in[0,\pi]$ es el ángulo entre ellos,
  \[
    \pp{u}{v}=\norma{u}\,\norma{v}\cos\theta .
  \]}

  \vspace{0.1cm}
  {\small \textit{Demostración} (ley de cosenos). El tercer lado del
  triángulo formado por $u$ y $v$ es $v-u$, de modo que
  \[
    \norma{v-u}^2=\norma{u}^2+\norma{v}^2
      -2\norma{u}\norma{v}\cos\theta .
  \]
  Por otra parte, desarrollando con las propiedades de la diapositiva
  anterior,
  \[
    \norma{v-u}^2=\pp{v-u}{v-u}
    =\norma{v}^2-2\pp{u}{v}+\norma{u}^2 .
  \]
  Igualando ambas expresiones y cancelando se obtiene el resultado.
  \hfill$\blacksquare$}

  \vspace{0.2cm}
  \destaca{\textbf{Consecuencia.} Perpendicularidad $\iff$ producto
  punto igual a cero. El giro con el que construimos el normal produjo un cero
  porque el ángulo era de $90^\circ$; ahora sabemos que no fue casualidad.}
\end{frame}
% ============================================================
\begin{frame}{La desigualdad del triángulo}
  En la semana 1 enunciamos sin demostrar la desigualdad del
  triángulo:
  \[
    d(P,Q)\ \le\ d(P,R)+d(R,Q) .
  \]
  Con $u=R-P$ y $v=Q-R$, basta demostrar
  $\norma{u+v}\le\norma{u}+\norma{v}$:
  \[
    \norma{u+v}^2=\norma{u}^2+2\,\pp{u}{v}+\norma{v}^2
    \;\le\;
    \norma{u}^2+2\,\norma{u}\norma{v}+\norma{v}^2 ,
  \]
  puesto que $\pp{u}{v}=\norma{u}\norma{v}\cos\theta$ y
  $\cos\theta\le 1$. \hfill$\blacksquare$

  \vspace{0.4cm}
  \destaca{La igualdad ocurre exactamente cuando $\cos\theta=1$, es
  decir, cuando $R$ pertenece al segmento $PQ$, tal como lo predijimos
  geométricamente en la semana 1. El pendiente que dejamos anotado en
  el pizarrón queda \textbf{demostrado}.}
\end{frame}
% ============================================================
\begin{frame}{Cierre de la semana}
  \textbf{Lo que vimos}
  {\small
  \begin{itemize}
    \item Puntos y vectores: qué operaciones tienen sentido y cuáles no
    \item La forma paramétrica $X=P+t\,d$, la pendiente, el giro de
      $90^\circ$, el normal y la ecuación general
    \item Las cinco formas de la recta y la mediatriz como lugar
      geométrico
    \item El producto punto, el teorema del ángulo, y la desigualdad
      del triángulo, ya demostrada
  \end{itemize}}

  \vspace{0.2cm}
  \textbf{Para la próxima semana}
  {\small
  \begin{itemize}
    \item Cuestionario de práctica de la semana 2 en Moodle
      (intentos ilimitados; no califica)
    \item Notas de la semana 2 en Moodle, al cierre de la semana
    \item Lehmann: serie de ejercicios sobre las formas de la ecuación
      de la recta
    \item Terminar la \textbf{Construcción 1} y entregarla en Moodle
      el domingo
  \end{itemize}}

  \vspace{0.2cm}
  \destaca{\textbf{Pregunta para la próxima sesión:} la ecuación
  $5x-2y=20$ contiene, a la vista, un vector perpendicular a la recta
  y sus dos cortes con los ejes. Identifíquenlos. Y el anuncio de la
  semana 3: con el producto punto mediremos el ángulo entre dos rectas
  y obtendremos \textbf{todas} las distancias del curso.}
\end{frame}
% ============================================================
%  Página de cierre del archivo.
%  Ajustar la URL cuando el libro quede publicado en GitHub Pages.
\begin{frame}[plain]
  \centering

  \vspace{1.0cm}
  \begin{tikzpicture}[scale=0.55,>=Stealth]
    \draw[very thin,color=gray!30] (-2.6,-1.4) grid (3.6,3.1);
    \draw[->,thick] (-2.6,0) -- (3.8,0);
    \draw[->,thick] (0,-1.4) -- (0,3.3);
    \draw[turquesaGA,very thick] (-2.4,-1.2) -- (3.4,1.7);
    \draw[terracota,very thick,->] (-0.4,-0.2) -- (1.6,0.8);
    \draw[dorado,very thick,->] (-0.4,-0.2) -- (-1.4,1.8);
  \end{tikzpicture}

  \vspace{0.5cm}

  {\Large\color{turquesaGA} ¿Preguntas?}\\[10pt]
  {\small\color{grisT}
  \url{https://rosaliahdez.github.io/notas-geometria-analitica/} \\
  Comunicación por Teams · Cuestionarios y entregas en Moodle (Ramanujan)}

  \vspace{0.7cm}

  {\tiny\color{grisT}
  Texto base: Ramírez Galarza, \textit{Geometría analítica. Una
  introducción a la geometría}, UNAM. ---
  Ejercicios: Lehmann, \textit{Geometría analítica}, Limusa.}
\end{frame}
\end{document}
